\Needspace{16\baselineskip}
\fancyhead[L]{\small VESTA. JUNO. PALLAS. CERES.}
\begin{center}
VESTA. JUNO. PALLAS. CERES.
\end{center}

\begin{center}
\rule{4cm}{0.4pt}\\
\fontsize{10}{12}\selectfont Göttingische gelehrte Anzeigen. Issue 107. Pp. 1065..1068. 4 July 1808.\\
\rule{4cm}{0.4pt}
\end{center}

On the night of 22 to 23 June, the new planet \emph{Vesta} was found again at our observatory and observed there again for the first time, after being invisible here for almost a whole year. It already looks exactly like an eighth-magnitude star, and its change of position became quite noticeable during the hour of observation up to daybreak. The mean of all the comparisons with an eighth-magnitude star, whose apparent position was reduced from the Hist. céleste and taken to be


\begin{center}\fontsize{10}{12}\selectfont\setlength{\tabcolsep}{3pt}\renewcommand{\arraystretch}{1.15}\begin{tabular}{l}
$352^\circ\;44'\;1''$ right ascension,\\
$8^\circ\;58'\;9''$ south declination
\end{tabular}\end{center}


\noindent gave the following result:


\begin{center}\fontsize{10}{12}\selectfont\setlength{\tabcolsep}{3pt}\renewcommand{\arraystretch}{1.15}\begin{tabular}{lrrr}
1808 & \shortstack{Mean time\\at\\Göttingen} & \shortstack{Apparent right\\ascension\\of Vesta} & \shortstack{South\\declination\\of Vesta}\\\hline
June 22 & $13^{\text{u}}\,46^{\text{m}}\,25^{\text{s}}$ & $353^\circ\;12'\;17''$ & $9^\circ\;11'\;55''$\\
\end{tabular}\end{center}


\par\indent The ephemeris calculated from Professor \textsc{Gauss}’s third elements and printed in the September 1807 issue of the Monatl. Correspondenz gives a right ascension $8'$ too small and a declination $3'$ too large. Those elements had already been calculated in July of the previous year using all the observations then available, the last of which was made at Lilienthal on 11 July. The Milan astronomers followed the planet much longer (until 26 Sept.), and their observations would subsequently have made a much more accurate determination of the orbit possible. However, although the later observations already differed from the third elements by more than one minute according to Professor \textsc{Harding}’s calculation, Professor \textsc{Gauss} considered a further correction of the elements unnecessary at the time. He also considered the ephemeris already calculated more than accurate enough to find the planet again, as the result has now confirmed. Professor \textsc{Gauss} later found a reason to recalculate the elements. His aim was to illustrate a section of his work on the orbits of celestial bodies, now in press, with a carefully chosen example. The most accurate final result possible from the previous year’s observations lay outside that aim. The calculation used one of the last Mi\SourceLeaf{309}lan observations, which, in keeping with the purpose of the calculation, had \emph{not} first been cleared of its probable observational error. Even so, the result could be expected to be considerably more accurate than the third elements. We therefore include it here, as it has yet to appear in print elsewhere:

\begin{center}
\emph{Elements of Vesta.}
\end{center}


\begin{center}\fontsize{9}{11}\selectfont\setlength{\tabcolsep}{3pt}\renewcommand{\arraystretch}{1.15}\begin{tabular}{lr}
Epoch 1807, Paris meridian .\;.\;.\;.\;.\;.\;.\;.\; & $168^\circ\;10'\;47''6$\\
Daily tropical mean motion .\;.\;.\;.\;.\;.\;.\;.\; & $978''8588$\\
Perihelion .\;.\;.\;.\;.\;.\;.\;.\;.\;.\;.\;.\;.\;.\;.\; & $249^\circ\;57'\;6''5$\\
Ascending node .\;.\;.\;.\;.\;.\;.\;.\;.\;,\;.\;.\; & $103\;11\;57$\\
\multicolumn{2}{l}{\quad both for 1807 and fixed with respect to the stars}\\
Inclination of the orbit .\;.\;.\;.\;.\;.\;.\;.\;.\;.\;.\;.\;.\; & $7\;8\;21$\\
Eccentricity .\;.\;.\;.\;.\;.\;.\;.\;.\;.\;,\;.\;.\;.\;.\;.\; & $0{,}0880158$\\
Logarithm of the mean distance .\;.\;.\;\hspace{.5em}.\;.\;.\;.\; & $0{,}3728980$\\
\end{tabular}\end{center}


\Needspace{9\baselineskip}
\par\indent The position calculated from \emph{these} elements for the time of the observation above is as follows:


\begin{center}\fontsize{10}{12}\selectfont\setlength{\tabcolsep}{3pt}\renewcommand{\arraystretch}{1.15}\begin{tabular}{rrrr}
right ascension & Difference & Declination & Difference\\\hline
$353^\circ\;14'\;30''$ & $+2'\;13''$ & $9^\circ\;12'\;26''$ & $+31''$\\
\end{tabular}\end{center}


\par\indent It should also be noted that the comparison star was very unfavorably placed for determining the declination, which is therefore rather unreliable. We believe we can vouch all the more for the right ascension, since the Hist. cél. gives two observations of the comparison star that yielded results in very good agreement. Professor \textsc{Gauss} took refraction into account in the reduction by a method of his own.

\par\indent \emph{Juno}, too, was found again here as early as 20 June. It is still very faint and looks only like a star of magnitude 10..11. The following observations have been possible so far: [They appear above on page 302 in the letter to \textsc{von Lindenau} dated 27 June 1808.]

\par\indent The position of the comparison star was determined from the Histoire céleste:


\begin{center}\fontsize{10}{12}\selectfont\setlength{\tabcolsep}{3pt}\renewcommand{\arraystretch}{1.15}\begin{tabular}{l}
right ascension $315^\circ\;24'\;58''$\\
south declination $2^\circ\;4'\;20''$
\end{tabular}\end{center}


\par\indent Accordingly, the seventh elements, published in issue 14 of these Anzeigen, differ from the observation by $11'$ in right ascension and $2'$ in declination. The ephemeris (Monatl. Corresp., February) gives the former too small and the latter too large. These differences can already be attributed in part to the effect of perturbations.

\par\indent \emph{Pallas} had already been found again on 20 May, but poor weather, moonlight and too great a distance from stars with well-determined positions have so far prevented observations. By a rough estimate, the ephemeris in the December issue of the Monatl. Corresp. seemed to give a good declination and a right ascension $2'$ too small.

\par\indent \emph{Ceres} has so far been too low in the morning twilight, but we hope to be able to report observations of it soon as well.
